\chapter{Woher die Startwerte kommen}\label{kap:startwerte}

Das Modell braucht für jede Mannschaft einen ELO-Wert zum Saisonstart. Er
steht in der Stammdatei \datei{RCode/TeamList_<Saison>.csv}, die als
gepflegtes Stammdatenblatt gilt (ADR~0007). Die Datei hat die Spalten
\code{TeamID}, \code{ShortText}, \code{Promotion}, \code{InitialELO},
\code{League}, \code{Region} und \code{Name}.

\section{Saisonwechsel}

Zum Saisonwechsel rechnet das Werkzeug
\datei{scripts/season_transition.R} die ELO am Ende der alten Saison aus.
Dazu ruft es denselben Endpunkt \code{POST /league-details} auf, den auch die
Website benutzt, je Liga mit deren Tormodell. Einen eigenen ELO-Walk in R
gibt es seit September 2026 nicht mehr (Issue~\#146). Der frühere hatte mit
einem Heimvorteil von 100 statt 40 gerechnet.

\begin{description}[font=\sffamily\bfseries\small]
  \item[Übertrag 1:1] Die Start-ELO der neuen Saison ist die End-ELO der
        alten, ohne Rückführung zur Mitte und unabhängig von Auf- oder
        Abstieg. Ein Aufsteiger nimmt seine Stärke mit.
  \item[Neue Mannschaften] Wer in keiner beobachteten Liga gespielt hat,
        bekommt einen Basiswert: in der 3.~Liga den Mittelwert der End-ELO der
        vier Absteiger der Vorsaison (Ersatzwert 1\,046), in der
        2.~Bundesliga 1\,350, in der Bundesliga 1\,500, sonst 1\,200
        (\datei{RCode/elo_aggregation.R:404--550}).
\end{description}

\section{Eichung für Ligen ohne Historie}\label{sec:eichung}

Mit der Erweiterung auf zehn Ligen im September 2026 kamen sieben Ligen ohne
eigene ELO-Historie dazu: die fünf Regionalligen und die beiden
Frauen-Ligen. Ihre Startwerte hat das einmalige Offline-Skript
\datei{scripts/calibrate_historical_elo.R} geeicht (ADR~0003):

\begin{enumerate}
  \item \textbf{Walk je Liga} über die verfügbare Historie, mit den
        Startwerten 1\,156 (3.~Liga), 920 (Regionalligen) und 1\,500 (Frauen).
        Die End-ELO einer Saison ist der Start der nächsten. Neue Mannschaften
        beginnen beim Startwert. Gewertet werden nur reguläre Spiele der
        Hauptrunde in chronologischer Reihenfolge.
  \item \textbf{Regionalligen an die 3.~Liga.} Der ganze Block wird
        verschoben, und zwar um die mittlere ELO-Differenz zwischen
        Drittliga-Absteigern und Regionalliga-Aufsteigern über die Wechsel
        2020--2025:
        \[
          \Delta_{RL} = \overline{E}(\text{Absteiger 3.~Liga}) - \overline{E}(\text{Aufsteiger RL}).
        \]
  \item \textbf{Staffeln gegeneinander.} Jede Staffel wird um die Hälfte der
        Abweichung verschoben, die ihre Aufsteiger in der 3.~Liga gegenüber
        dem Mittel aller Staffeln zeigen. Die Dämpfung auf die Hälfte
        berücksichtigt, dass je Staffel nur zwei bis sechs Aufsteiger
        vorliegen.
  \item \textbf{2.~Frauen-Bundesliga an die Frauen-Bundesliga}, analog zu
        Schritt~2.
  \item \textbf{Frauen an die Herren}, per Konvention: Die Frauen-Ligen werden
        so verschoben, dass der Mittelwert der Frauen-Bundesliga dem der
        Herren-Bundesliga entspricht (1\,682). Das ist ausdrücklich keine
        Messung, denn zwischen beiden findet kein Austausch statt.
\end{enumerate}

Alle Verschiebungen sind additiv, die Abstände innerhalb einer Liga bleiben
erhalten. Weil der Walk die Summe der ELO erhält, ist keine Iteration nötig.

\section{Ein Tormodell je Wechselgemeinschaft}\label{sec:wechselgemeinschaft}

Die ELO ist das Einzige, was eine Mannschaft über eine Ligagrenze mitnimmt.
Nur wenn Ligen, die Mannschaften austauschen, dasselbe Tormodell benutzen,
bedeutet ein ELO-Wert diesseits und jenseits der Grenze dasselbe. Maßgeblich
ist deshalb nicht die einzelne Liga, sondern die \emph{Wechselgemeinschaft}
(ADR~0004):

\begin{itemize}
  \item \textbf{Herren} (Ligen 78, 79, 80, 83--87) tauschen untereinander
        Mannschaften aus und teilen ein Tormodell.
  \item \textbf{Frauen} (82, 1034) sind untereinander verbunden, nie mit den
        Herren, und haben daher ein eigenes.
\end{itemize}

Die Frauen-Werte $s_F, c_F$ sind per Maximum-Likelihood an 1\,917 Spielen
geschätzt und außerhalb der Stichprobe validiert. Der Brier-Score sinkt um
2{,}96~\%, der Log-Loss um 2{,}66~\%, und das neue Modell ist in allen fünf
Testsaisons besser
(\datei{docs/reports/2026-09-05-frauen-tormodell.md}). Ein eigener
Heimvorteil für die Frauen hat sich nicht als signifikant erwiesen
(Schätzwert 28, $p = 0{,}23$). Es bleibt bei 40.

\subsection*{Wie der Heimvorteil zu 40 kam}

Bis September 2026 rechnete das Modell mit 65 ELO-Punkten. Die
Neukalibrierung integriert über eine realistische Verteilung von
ELO-Unterschieden und trifft mit 40 die gemessenen Anteile der Bundesliga
2024 und 2025:

\begin{center}
\begin{tabular}{@{}lccc@{}}
\toprule
& Heimsieg & Remis & Auswärtssieg \\
\midrule
Modell mit $H = 40$ & 41{,}0\,\% & 24{,}4\,\% & 34{,}7\,\% \\
gemessen            & 41{,}2\,\% & 25{,}0\,\% & 33{,}8\,\% \\
\bottomrule
\end{tabular}
\end{center}

Auf die Saisonprognosen wirkt die Änderung mit weniger als einem
Prozentpunkt (ADR~0002, Nachtrag;
\datei{docs/reports/2026-09-05-heimvorteil-vorher-nachher.md}).
